\scanpage{4}
\[
\begin{aligned}
\text{b)}\ & h_i\quad(0\leq i\leq n)\quad \text{engendrent } C\quad\text{i.e. lin. indép.}\\
\text{c)}\ & h_i\notin H_{i+1}\quad(0\leq i\leq n-1)\quad\text{en tte fibre}
\end{aligned}
\]

\margin{non autodual sauf si $\rho$ inv.}
$\widetilde C^\infty_{vect}$ variante vectorielle $(V,\varphi_{\unsure{G|g|\mathfrak G}}^V,L_0)$, $V$ fibré vect. rg $n+1$, $\varphi_G^V:\unsure{\mathscr G|G|\mathfrak G}\longrightarrow\operatorname{Aut}(V)$, $L_0$ sous-fibré vect. de rg $1$, avec

a) les $\tau_i^V$ $0\leq i\leq n$ ps. réflexions, strictes pour $1\leq i\leq n$, et sub. à $L_0$ si $i=0$ ($\Rightarrow L_i,H_i,L'_i$ si $1\leq i\leq n$)
\[
\begin{aligned}
\text{b)}\ & V=\bigoplus_{0\leq i\leq n} L_i\\
\text{c)}\ & h_i\notin H_{i+1}\qquad(0\leq i\leq n-1)
\end{aligned}
\]
Alors $L_0\simeq L_1^\vee\simeq L_1\simeq L_2^\vee\simeq L_2\ldots\simeq L_n^\vee\simeq L_n$, soit $L$.

On appelle \uline{rigidification} la donnée d'une base de $L$, (p. ex. d'une base $a_0$ de $L_0$, d'où $a_i$ de $L_i$, $a'_i$ de $L'_i$).

$C^\infty_{vect}$ cat. des objets de $\widetilde C^{\unsure{n|\infty}}_{vect}$ rigidifiés

\margin{$n\geqslant 2$, les $\mu_i=-2$}
\margin{non autodual sauf si $\rho$ inv.}
$C^\infty_{pr.\,orth}$ cat des objets $(C,Q,\varphi_G^{C,Q},b_0)$, où $C$, $\deleted{\unsure{\varphi_G^C|\varphi_{\mathfrak G}^C}}$ fibré projectif rg $n$, $Q$ hyper\inserted{$\unsure{\infty}$}quadrique \deleted{la donnée \unsure{revient} par $C^\infty_{prog}$, et $Q$ hyperquadrique inv. de $C$,}
\[
\varphi_G^{C,Q}:\mathscr G_n\longrightarrow\operatorname{Aut}(C,Q),
\]
$b_0$ section de $C$, tels que

a) les $\tau_i^{C,Q}$ $(0\leq i\leq n)$ réfl. orth, strictes pour $1\leq i\leq n$, sub. à $h_0$ si $i=0$ ($\Rightarrow h_i,H_i\quad1\leq i\leq n$)
\[
\begin{aligned}
\text{b)}\ &h_i\quad (0\leq i\leq n)\quad\text{engendrent }C\\
\text{c)}\ &h_i\notin H_{i+1}\quad\text{en tte fibre}\quad(0\leq i\leq n-1)
\end{aligned}
\]

\margin{\uline{NB} \deleted{Dans} le cas \deleted{\unsure{propre}} \inserted{$\beta_0\beta_1$ inv.} (plus \inserted{gén} $\beta_0$ \deleted{\unsure{prem.}} inv., \uline{et} $\deleted{(\beta_1,2)\mathcal O_S=\mathcal O_S}$ \deleted{($\beta_1$ inv. en tte \unsure{pt} de car. $2$)} $(\beta_1,\deleted{\unsure{\beta}}4-\beta_0)\mathcal O_S=\mathcal O_S$ i.e. $\beta_1=\alpha_1+2$ et $4-\beta_0=2-\alpha_0$ non nuls simultanément) b) équivaut au \uline{fait} que sur aucune fibre, il n'y a d'hyperplan \uline{de $C$} invariant par $G$ \ldots\ (cette remarque valant aussi pour $C^\infty_{vect}$, $C^\infty_{pr}$ \ldots. expliciter}
\uline{ReNB} si $\tau_i^C,\tau_{i+1}^C$ réflexions \uline{str.} orth. i.e. $\beta_0\ldots\beta_i$ inv. Alors $h_i\notin H_{i+1}$ en tte fibre équivaut à $\underbrace{\tau_i\tau_{i+1}}_{\rho_i}\ne\tau_{i+1}\tau_i$ \inserted{tte fibre} i.e. $\rho_i^2\ne\mathrm{id}$ en tte fibre. Donc si $\beta$ inv. c) se remplace par ces conditions pour $0\leq i\leq n-1$ \ldots

% archive sheet number: 3

